\section{The tree coefficient \texorpdfstring{$A_g$}{A(g)}}\label{sec:amplitude}

Throughout this section $P$ is a polygon satisfying (G1); (G2) is not used.
The object is the finite sum of the main text (\mt{eq:amplitudevertex}--\mt{eq:amplitude}), written
here as a recursion over interval compositions, which is the form in which
every law is proved.

\subsection{Roots, boundary words, gates}

\begin{definition}[root and boundary word]\label{def:root}
A \emph{root} of a polygon is one of its edges; on a representative it is an
edge index $g\in\ZZ/n$, the root edge being $E_g$ from $\mu_g$ to
$\mu_{g+1}$, and the shift carries the root $g$ of $P$ to the root $g-1$ of
$\sigma P$. Everything in this section is a function of the pair (polygon,
root). The \emph{boundary word} of $(P,g)$ is the sequence of
vertices
\[
a_0=\mu_{g+1},\ a_1=\mu_{g+2},\ \ldots,\ a_{N}=\mu_{g+n}=\mu_g,\qquad N=n-1,
\]
so that consecutive boundary vertices $a_{k-1},a_k$ are joined by the
\emph{leaf} edge $E_{g+k}$, $1\leq k\leq N$; the leaves are the $n-1$ edges
other than the root, in traversal order. An \emph{interval} is a pair
$[i,j]$ of integers $0\leq i<j\leq N$; it has $j-i$ leaves. A
\emph{composition} of $[i,j]$ is a sequence $\pi=(r_0<r_1<\cdots<r_s)$ with
$r_0=i$, $r_s=j$, $s\geq1$; its \emph{parts} are the intervals
$[r_{k-1},r_k]$, and $|\pi|=s$.
\end{definition}

\begin{definition}[gates]\label{def:gates}
Let $\pi=(r_0<\cdots<r_s)$ be a composition of $[i,j]$. At each interior cut
$r_k$, $1\leq k\leq s-1$, put
\[
d_k(\pi)=\chi\bigl(a_{r_{k+1}},a_{r_k},a_{r_{k-1}}\bigr),\qquad
h_k(\pi)=\chi\bigl(a_j,a_{r_k},a_i\bigr),
\]
where $\chi(x,y,z)=\sgn\det(y-x,z-x)$ for points $x,y,z$; $d_k$ is the
\emph{near sign} and $h_k$ the \emph{far sign}. The \emph{ordinary weight}
and the \emph{root weight} of $\pi$ are
\[
V^+(\pi)=\prod_{k=1}^{s-1}d_k\,\Theta\bigl(-h_kd_k\bigr),\qquad
V^-(\pi)=\prod_{k=1}^{s-1}d_k\,\Theta\bigl(h_kd_k\bigr),
\]
with empty products equal to $1$ (so $V^\pm(\pi)=1$ when $s=1$).
\end{definition}

\begin{lemma}[gates are defined on (G1)]\label{lem:gates-nonzero}
If $P$ satisfies (G1) then every near and far sign is $\pm1$ and every
$\Theta$ in Definition~\ref{def:gates} is evaluated at $\pm1$. On that
domain
\[
d\,\Theta(-hd)=\frac{d-h}{2},\qquad d\,\Theta(hd)=\frac{d+h}{2}\qquad(d,h\in\{\pm1\}),
\]
and $V^+(\pi)=0$ for every composition with exactly two parts.
\status{proved (refereed: bench A 2026-09-12T06:44:52Z, bench B 2026-09-12T06:57:24Z; SM1-carried, read on SM12, no text change, re-read on SM13)}
\end{lemma}
\begin{proof}
The three boundary vertices entering a gate are pairwise distinct vertices of
$P$ (the cut positions are distinct and $a_0\neq a_N$ by (G1)), so the signs
are nonzero. The two identities are checked on the four sign pairs. For
$s=2$ the only cut is $r_1$ and $d_1=\chi(a_j,a_{r_1},a_i)=h_1$, so
$\Theta(-h_1d_1)=\Theta(-1)=0$.
\end{proof}

\subsection{The coefficient}

\begin{definition}[open sums and the tree coefficient]\label{def:treesum}
Let $P$ satisfy (G1) and fix a root $g$. Define $b_{[i,j]}\in\ZZ$ for every
interval by recursion on the number of leaves:
\[
b_{[i,i+1]}=1,\qquad
b_{[i,j]}=-\sum_{\substack{\pi\text{ composition of }[i,j]\\|\pi|\geq2}}
V^+(\pi)\prod_{[r_{k-1},r_k]\in\pi}b_{[r_{k-1},r_k]}\quad(j-i\geq2).
\]
The \emph{tree coefficient} of $P$ at the root $g$ is
\[
A_g(P)=\sum_{\pi\text{ composition of }[0,N]}
V^-(\pi)\prod_{[r_{k-1},r_k]\in\pi}b_{[r_{k-1},r_k]},
\]
the sum including the one-part composition, whose term is $b_{[0,N]}$. The main
text's $A(P)$ is $A_n(P)$, rooted at the last edge $E_n$ from $\mu_n$ to
$\mu_1$.
\end{definition}

\begin{lemma}[equivalence with the plane-tree sum]\label{lem:treesum-trees}
Let $\mathfrak T_g(P)$ be the set of rooted plane trees whose leaves are the
$n-1$ leaf edges in traversal order, whose root vertex has one or more
children, and whose other internal vertices (the \emph{ordinary} vertices)
have at least two children each; every internal vertex is labelled by the
composition of its leaf interval into the leaf intervals of its children.
Then
\[
A_g(P)=\sum_{T\in\mathfrak T_g(P)}(-1)^{N_+(T)}\,
V^-(\pi_{\rm root})\prod_{v\text{ ordinary}}V^+(\pi_v),
\]
where $N_+(T)$ is the number of ordinary vertices. The identity is formal:
the bijection of the proof matches the terms of the recursion of
Definition~\ref{def:treesum} with the trees term by term, so it holds with
the gate values $V^\pm(\pi)$ replaced by independent formal variables, as an
identity in the polynomial ring over $\ZZ$ in those variables, and
specializes to every evaluation of them. This is \mt{eq:amplitude},
whose vertex weights are Definition~\ref{def:gates} with
$m=s$, $a_k$ the boundary vertices of the cuts, and whose sign
$(-1)^{N_+}$ is the one minus per ordinary vertex. \status{new (round~5: the formal-variable clause exported --- the grafting bijection of the proof uses nothing about the gate values, so the identity holds in the polynomial ring in those values and specializes to every evaluation, F-25-150; the remainder is the SM1 proof, audited there)}
\end{lemma}
\begin{proof}
Removing the top vertex of an ordinary tree on $[i,j]$ leaves one subtree per
part of its composition; grafting is the inverse. The minus sign of the top
vertex is the minus sign in the recursion for $b$. The root vertex carries
$V^-$, no sign, and may have one child, which is the one-part composition.
\end{proof}

\begin{remark}
Ordinary binary vertices contribute $0$ (Lemma~\ref{lem:gates-nonzero}),
which is why the main text's Berends--Giele form, \mt{eq:BGG}, sums over $m\geq3$.
The main text's ``$V^+$ propagators are discarded'' is the requirement that
ordinary vertices have at least two children.
\end{remark}

\subsection{Chamber constancy and the resummation}

\begin{proposition}[chirotope chambers]\label{prop:A-chamber}
If two polygons satisfying (G1) have the same chirotope, then for every root
$g$ every composition weight, every $b_{[i,j]}$ and $A_g$ agree. In
particular $A_g$ is constant on every labelled chamber (the root fixed), and
constant along any path of (G1)-polygons with constant chirotope. \status{proved (refereed: bench A 2026-09-12T06:44:52Z, bench B 2026-09-12T06:57:24Z; SM1-carried, read on SM12, repaired in round 12 (F-25-191), re-read on SM13)}
\end{proposition}
\begin{proof}
Every quantity in Definitions~\ref{def:gates} and~\ref{def:treesum} is a
function of the signs $\chi(a_x,a_y,a_z)$, which are chirotope entries. A
labelled chamber is a component of $\mathcal U_n$, an open set
(Proposition~\ref{prop:chambers}), hence path connected, and the chirotope
is constant along every path in $\mathcal U_n$
(Proposition~\ref{prop:chambers}); so $A_g$ is constant on each labelled
chamber.
\end{proof}

\begin{definition}[the near--far arrays]\label{def:nearfar}
Over a commutative ring in which $2$ is invertible, let $D(x,y,z)$ and
$H(x,y,z)$ be arbitrary scalars indexed by increasing triples
$0\leq x<y<z\leq N$. For an array $X=(X_{ij})_{0\leq i<j\leq N}$ put
\[
\Phi_{D,H}(X)_{ij}=\sum_{\pi=(i=r_0<\cdots<r_s=j)}
\Bigl(\prod_{k=1}^{s-1}\frac{D(r_{k-1},r_k,r_{k+1})-H(i,r_k,j)}2\Bigr)
\prod_{k=1}^sX_{r_{k-1}r_k},
\]
$F_H=\Phi_{0,H}$, $G_D=\Phi_{D,0}$, and $E_{ij}=1$ if $j=i+1$ and $0$
otherwise. The \emph{geometric arrays} are $D(x,y,z)=\chi(a_z,a_y,a_x)$ and
$H(x,y,z)=\chi(a_z,a_y,a_x)$ with the boundary vertices of $(P,g)$.
\end{definition}

\begin{lemma}[factorization and far-only output]\label{lem:farout}
\begin{enumerate}
\item[(i)] Over any commutative ring in which $2$ is invertible,
$\Phi_{D,H}$ has a two-sided polynomial inverse and
\begin{equation}\label{afr:factor}
\Phi_{D,H}=F_H\circ G_D.
\end{equation}
\item[(ii)] With the geometric arrays, $\Phi_{D,H}(b)=E$ and
$A_g(P)=\Phi_{D,-H}(b)_{0N}$. Consequently, with $c=F_H^{-1}(E)$,
\begin{equation}\label{afr:far-output}
A_g(P)=F_{-H}(c)_{0N}.
\end{equation}
The complete output is independent of $D$ when $H$ is held fixed.
\item[(iii)] For an interval $J=[i,j]$ put
$\mathcal E_J=F_H(c)_J$ and $\mathcal B_J=F_{-H}(c)_J$.
Then $\mathcal E_J=0$ if $j-i\geq2$, and
$\mathcal E_J=\mathcal B_J=1$ if $j-i=1$.
For $j-i\geq2$, if $(a_i,\ldots,a_j)$ satisfies (G1),
$\mathcal B_J$ is its tree coefficient rooted at its closing edge
from $a_j$ to $a_i$.
\end{enumerate}
\status{proved (refereed: bench A, 2026-09-05T15:51:04Z; transcribed from RC mt:factorization, mt:far-only (C022))}
\end{lemma}
\begin{proof}
For (i), the one-part composition in coordinate $[i,j]$ contributes
$X_{ij}$ with coefficient one. Every other composition has at least two
parts, all of shorter interval length. Given any target $Y$, first put
$X_{i,i+1}=Y_{i,i+1}$ and then, in increasing interval length, subtract
the already determined nonunary sum from $Y_{ij}$. This uniquely solves
$\Phi_{D,H}(X)=Y$, by polynomial operations in the finite arrays.
Writing the solution map as $\Psi$, construction gives
$\Phi_{D,H}\circ\Psi=\mathrm{id}$. Applying uniqueness to the target
$\Phi_{D,H}(X)$ gives $\Psi\circ\Phi_{D,H}(X)=X$, the other inverse.

Expand $F_H(G_D(X))_{ij}$. Its terms are indexed by an outer
composition of $[i,j]$ and an inner composition of each outer part.
Their ordered union is a refined composition
$\rho=(r_0<\cdots<r_s)$, with the subset of its interior cuts marked
as outer cuts. Conversely, cut $\rho$ at the marked cuts to recover
the outer parts and restrict $\rho$ to each part to recover its inner
composition. This is a bijection, including empty and full marked
subsets and unary inner or outer compositions.

At a marked cut $r_k$ the coefficient is $-H(i,r_k,j)/2$.
At an unmarked cut its immediate neighbors in $\rho$ are precisely
its neighbors inside the same inner part, since that part's endpoints
are also in $\rho$. Its coefficient is
$D(r_{k-1},r_k,r_{k+1})/2$. Each interior cut supplies one factor,
and the product of child coordinates is exactly the refined product.
Summing over every subset of marked cuts independently chooses the
near or negative far summand at each cut. Distributivity yields the
product of their sums, precisely the coefficient of $\rho$ in
Definition~\ref{def:nearfar}. This proves \eqref{afr:factor}, with all
neighboring near gates accounted for.

For (ii), the identities $d\Theta(-hd)=(d-h)/2$ and
$d\Theta(hd)=(d+h)/2$, for $d,h\in\{\pm1\}$, turn the open-sum
recursion into $\Phi_{D,H}(b)=E$ and the full root sum into
$\Phi_{D,-H}(b)_{0N}$. In particular its unary root term is retained.
By \eqref{afr:factor}, $F_H(G_D(b))=E$; uniqueness of the inverse
gives $G_D(b)=c$. Applying the factorization with $-H$ gives
\begin{equation}\label{afr:near-freedom}
\Phi_{D,-H}(b)=F_{-H}(G_D(b))=F_{-H}(c).
\end{equation}
This proves the assertion about the complete output. It asserts no
near-array independence for individual trees or open sums.

For (iii), restriction to $J$ commutes with all coordinate equations:
each equation on a subinterval of $J$ reads only its own subintervals
and triples. Triangular uniqueness therefore identifies the restriction
of $c$ with the inverse constructed directly on that restricted word.
The equation $F_H(c)=E$ gives the claimed $\mathcal E$ values.
On a one-leaf interval the only composition is unary, its inverse
coordinate is one, and both transforms equal one. This is a formal
open-word convention, not a two-gon amplitude.
On an interval with at least two leaves, the closed endpoint polygon
has at least three vertices. Cutting its closing edge from $a_j$ to
$a_i$ gives exactly the word $a_i,\ldots,a_j$. If it satisfies (G1),
part (ii) on that word identifies $\mathcal B_J$ with its amplitude
at that specified physical root. The auxiliary coordinates $c$ may be
rational and their binary coordinates need not vanish.
\end{proof}

\begin{theorem}[single-triple wall response]\label{thm:single-triple}
Let $t\mapsto P(t)$ be a wall germ with $Z_{\rm c}=\varnothing$,
$Z_{\rm pt}=\{K\}$ a single triple of pairwise distinct vertices,
and $\chi_K$ changing sign at zero. This includes types (F), (K),
(V), (E), and (C). Cut at a fixed physical root $g$ and let
$x<y<z$ be the critical boundary positions. Put
$I_*=[x,z]$, $L=[x,y]$, $R=[y,z]$, $F=[0,N]$, and
\begin{equation}\label{afr:wall-data}
\delta=\frac{H^+(x,y,z)-H^-(x,y,z)}2\in\{-1,1\}.
\end{equation}
Write the three wall points as $a_v=p_*+\xi_v\omega$ with
$\omega\ne0$ and distinct $\xi_v$. Set
\begin{equation}\label{afr:wall-epsilons}
\epsilon_L=\sgn\frac{\xi_y-\xi_x}{\xi_z-\xi_x},\qquad
\epsilon_R=\sgn\frac{\xi_z-\xi_y}{\xi_z-\xi_x}.
\end{equation}
For $J=L,R$ use the unchanged interval values of
Lemma~\ref{lem:farout}, and put
\begin{equation}\label{afr:wall-uv}
\mathcal U_J=\begin{cases}\mathcal E_J,&\epsilon_J=1,\\
\mathcal B_J,&\epsilon_J=-1,\end{cases}\qquad
\mathcal V_J=\begin{cases}\mathcal B_J,&\epsilon_J=1,\\
\mathcal E_J,&\epsilon_J=-1.\end{cases}
\end{equation}
If $I_*\ne F$, let $Q$ be the polygon obtained at $t=0$ by
replacing the boundary arc $a_x,\ldots,a_z$ with the single edge
from $a_x$ to $a_z$. Then $Q$ satisfies (G1), has at least three
vertices, retains the physical root $g$, and
\begin{equation}\label{afr:wall-proper}
A_g(P_+)-A_g(P_-)=\delta\mathcal U_L\mathcal U_R A_g(Q).
\end{equation}
If $I_*=F$, instead
\begin{equation}\label{afr:wall-full}
A_g(P_+)-A_g(P_-)
=\delta(\mathcal U_L\mathcal U_R+\mathcal V_L\mathcal V_R).
\end{equation}
Here $P_\pm$ are on sufficiently close punctured sides of the germ;
all gap and contracted coefficients are evaluated at their nondegenerate
wall data, equivalently at the same signs on either sufficiently close side.
\status{proved (refereed: bench A, 2026-09-05T15:51:04Z and 2026-09-05T19:42:49Z; transcribed from RC mt:wall-response (C022))}
\end{theorem}
\begin{proof}
The wall hypothesis makes every distinct-triple determinant except the
critical one nonzero at zero; after a common shrink its sign is the
same on both punctured sides. Neither gap contains the entire critical
triple, so both gap arrays and every nontrivial closed gap polygon are
nondegenerate and unchanged. Singleton gap values have only their formal
open-word meaning. At least one epsilon is positive: otherwise both
summands in $\xi_z-\xi_x=(\xi_y-\xi_x)+(\xi_z-\xi_y)$ would have
sign opposite to their nonzero sum.

Use the far-only representation of Lemma~\ref{lem:farout} on each side.
The changing triple appears as far data in exactly one coordinate formula
of $F_H$, namely $[x,z]$.  Induction in interval length therefore shows
that $c$ is unchanged on every interval not containing $I_*$.

In the $[x,z]$ equation every proper child is unchanged.  A composition
has a changing coefficient exactly when it cuts at $y$.  Splitting its
cut list at $y$ is a bijection with a pair of arbitrary compositions of
$L$ and $R$; concatenating the two lists is its inverse.  Its one
changing far gate $-H(x,y,z)/2$ has difference $-\delta$.
For every other cut in these compositions, the wall geometry gives
\begin{align}
 H(x,r,z)&=\epsilon_L H(x,r,y) &&(x<r<y),\label{afr:wall-left-sign}\\
 H(x,r,z)&=\epsilon_R H(y,r,z) &&(y<r<z).\label{afr:wall-right-sign}
\end{align}
For the first identity, the vectors from $a_x$ to $a_z$ and to $a_y$
differ by the scalar $(\xi_z-\xi_x)/(\xi_y-\xi_x)$.
Taking their determinants with $a_r-a_x$ gives the indicated sign ratio.
For the second use the vectors from $a_z$ to $a_x$ and to $a_y$;
their scalar ratio has sign $\epsilon_R$.  All these determinants are
nonzero at the wall.  Their signs on both sufficiently close sides are
therefore the wall signs, proving both displayed identities on those
sides without evaluating a zero gate.

Consequently the composition sum on a gap is $F_H(c)_J=\mathcal E_J$ if
$\epsilon_J=1$, and $F_{-H}(c)_J=\mathcal B_J$ if $\epsilon_J=-1$.
The changed nonunary sum in $F_H(c)_{xz}=E_{xz}=0$ is thus
$-\delta \mathcal U_L\mathcal U_R$.  Since the unary coefficient is $1$, moving this
sum to the other side gives
\begin{equation}\label{afr:wall-source}
 c^+_{xz}-c^-_{xz}=s_*,
 \qquad s_*=\delta \mathcal U_L\mathcal U_R.
\end{equation}
The nonunary part of $F_{-H}(c)_{xz}$ must be calculated separately.
Its changing gate is $+H(x,y,z)/2$, with difference $+\delta$;
all proper children are still unchanged.  Reversing the remaining far
signs exchanges $\mathcal E_J$ with $\mathcal B_J$ in each of the two gap sums.
Its direct coefficient change is therefore $+\delta \mathcal V_L\mathcal V_R$.

It remains to propagate \eqref{afr:wall-source} through larger intervals.
Contract $I_*$ in the boundary word to the formal leaf $e_*$, and let
$c^Q$ be the inverse $F_{H^Q}^{-1}(E)$ on that contracted word.
This inverse is defined even when the contracted word has only one
leaf; that is an open-word convention, not a two-gon amplitude.
For an interval containing $I_*$ write $I/e_*$ for its contraction.
We prove by interval length that
\begin{equation}\label{afr:wall-propagation}
 c^+_I-c^-_I=s_*c^Q_{I/e_*}\quad(I\supseteq I_*),
 \qquad c^+_I-c^-_I=0\quad(I\not\supseteq I_*).
\end{equation}
The smallest containing interval is $I_*$ itself.  Its contracted
coordinate is the leaf value $1$, so this is \eqref{afr:wall-source}.
For a strictly larger $I$, its far coefficient formula is unchanged.
In a top composition at most one child can contain the positive-length
interval $I_*$, because distinct children have disjoint interiors.
If no child contains it, every child coordinate is unchanged.  Otherwise
the product difference is exactly the difference of that unique child
times all other, unchanged child coordinates.  There is no product of
two changing factors or choice of side for the others.

The compositions having such a child are in bijection with compositions
of $I/e_*$.  Forward, erase the interior vertices of $I_*$ inside that
one child; no top cut lies there.  Backward, the distinguished leaf
$e_*$ belongs to one unique contracted child, and expanding it
recovers the original child and top cut list.  The operations are
inverse.  All top endpoints and cut vertices survive, so every far
gate has the same indices and signs.  Every other child is disjoint
from the interior of $I_*$, and its $c$ coordinate agrees with $c^Q$
by restriction and uniqueness of the unchanged triangular equations.
Insert the induction hypothesis on the changing child.  The difference
of the equation $F_H(c)_I=0$ becomes $s_*$ times the equation
$F_{H^Q}(c^Q)_{I/e_*}=0$.  The contracted interval has at least two
leaves when $I$ strictly contains $I_*$, so this latter right-hand
side is indeed zero.  Equivalently, its unary term is the negative
of its full nonunary sum, giving precisely
$c^+_I-c^-_I=s_*c^Q_{I/e_*}$.  This proves the induction with its sign.

When $I_*\ne F$, the same composition bijection, now including the
unary term, applies to the barred transform $F_{-H}$ on $F$.
Its coefficients are unchanged because its endpoints are not both $x,z$.
Substitution of \eqref{afr:wall-propagation} gives
$s_*F_{-H^Q}(c^Q)_{F/e_*}$.  At least one further leaf survives outside
$I_*$, so $Q$ has at least three vertices.  Its vertices have no zero
triple, since the only critical triple lost its middle occurrence.
Its endpoints and hence physical root survive.  Lemma~\ref{lem:farout}
identifies the last output with $s_*A_g(Q)$, proving
\eqref{afr:wall-proper}.

When $I_*=F$, all proper children are unchanged and the unary
contribution is \eqref{afr:wall-source}.  Add the separately computed
direct barred coefficient change $+\delta \mathcal V_L\mathcal V_R$ to obtain
\eqref{afr:wall-full}.  In this case there is no closed contracted
amplitude of arity at least three.
\end{proof}




\subsection{The laws}

\begin{definition}[induced roots]\label{def:induced-roots}
At a simple flat or cusp wall at $j$, the \emph{deletion map} on roots is
\[
D_j(g)=\begin{cases}\text{the fused edge }[\mu_{j-1},\mu_{j+1}]\text{ of }P\setminus j,& g\in\{j-1,j\},\\
\text{the edge of }P\setminus j\text{ with the same endpoints as }E_g,&\text{otherwise.}\end{cases}
\]
At a simple vertex--edge wall at $(M;a)$, with $d_1=[\mu_a,\mu_M]$ the
closing edge of $\lambda_1$ and $d_2=[\mu_M,\mu_{a+1}]$ the opening edge of
$\lambda_2$, the \emph{half map} on roots is
\[
H(g)=\begin{cases}(d_1,d_2),&g=a,\\
(E_g\text{ in }\lambda_1,\ d_2),&g\in\{M,M+1,\ldots,a-1\},\\
(d_1,\ E_g\text{ in }\lambda_2),&g\in\{a+1,\ldots,M-1\}.\end{cases}
\]
\end{definition}

\begin{theorem}[flat law for $A_g$]\label{thm:A-S3}
At a simple flat wall at $j$ with $n\geq4$, for every root $g$,
\begin{equation}\label{afr:S3}
A_g(P_{\rm right})-A_g(P_{\rm left})
=A_{D_j(g)}(P(0)\setminus j),
\end{equation}
where $P_{\rm right}$ is the side with $\tau_j=-1$, and
$D_j$ is the physical deletion map of Definition~\ref{def:induced-roots}.
\status{proved (refereed: bench A, 2026-09-05T15:51:04Z; transcribed from RC mt:s3 (C022))}
\end{theorem}
\begin{proof}
Put $Q=P(0)\setminus j$. Its vertices are a subset of the wall
vertices missing the middle occurrence of the sole zero triple. Hence
all its distinct-triple determinants are nonzero, so it satisfies (G1),
and $n-1\geq3$. The flat hypothesis puts the deleted vertex strictly
between its neighbors. Let $P^\pm$ denote the sides on which
$\chi(\mu_{j-1},\mu_j,\mu_{j+1})=\pm1$, and write
$u=(\mu_{j-1},\mu_j)$, $v=(\mu_j,\mu_{j+1})$,
$w=(\mu_{j-1},\mu_{j+1})$. The determinant defining this chirotope
is $\det(\widetilde\lambda_{j-1},\widetilde\lambda_j)$, since
$\mu_{j+1}-\mu_{j-1}=\widetilde\lambda_{j-1}+\widetilde\lambda_j$.
Thus it is $\tau_j$, and $P^-=P_{\rm right}$,
$P^+=P_{\rm left}$.

Write $A=\mu_{j-1}$, $M=\mu_j$, $B=\mu_{j+1}$.  In any root cut the
three critical occurrences have a cyclic permutation of the order
$A,M,B$.  A cyclic permutation is even, so their reverse-ordered far
sign is always $-\chi(A,M,B)$.  Thus the $\delta$ of
\eqref{afr:wall-data}, for $P^+$ minus $P^-$, is $-1$.

If $g$ is neither $u$ nor $v$, the cut order is $A,M,B$, with no
root seam between these consecutive edges.  Both gaps are singleton
leaves, and strict betweenness gives
$(\epsilon_L,\epsilon_R)=(1,1)$.  Hence $\mathcal U_L\mathcal U_R=1$.
Their span is proper: if both extreme vertices were the root endpoints,
the remaining root edge would close a three-vertex polygon, contrary
to $n\geq4$.  Contracting the two-edge span deletes exactly $M$
and retains $g$.  Formula~\eqref{afr:wall-proper} gives
$A_g(P^+)-A_g(P^-)=-A_g(Q)$.

If $g=u$, the cut order is $M,B,\ldots,A$, so the critical span is
full.  The gap $L=(M,B)$ is singleton and
$(\epsilon_L,\epsilon_R)=(-1,1)$: with $A=0$, $B=1$,
$M=t\in(0,1)$, the first ratio has numerator $1-t>0$ and
denominator $-t<0$, and the second has numerator $-1<0$.
The other gap $R$ has $n-2\geq2$ leaves; its closing root is
$w=(A,B)$ and its closed polygon is $Q$.  Therefore
$\mathcal U_L\mathcal U_R=0$ and $\mathcal V_L\mathcal V_R=A_w(Q)$.
Formula~\eqref{afr:wall-full} gives the desired negative of $A_w(Q)$.

If $g=v$, the cut order is $B,\ldots,A,M$.  Now $L$ is the
nontrivial deletion word with closing root $w$, $R=(A,M)$ is
singleton, and $(\epsilon_L,\epsilon_R)=(1,-1)$.
Indeed their numerators are $-1$ and $t$ with common denominator
$t-1<0$.  Again $\mathcal U_L\mathcal U_R=0$ and $\mathcal V_L\mathcal V_R=A_w(Q)$.
These three cases exhaust the physical roots and prove
\eqref{afr:S3} with precisely Definition~\ref{def:induced-roots}.
\end{proof}

\begin{theorem}[cusp law for $A_g$]\label{thm:A-S4}
At a simple cusp wall at $j$ with $n\geq4$ whose deletion
$Q=P(0)\setminus j$ satisfies (G1), for every root $g$,
\begin{equation}\label{afr:S4}
A_g(P_{\rm loop})-A_g(P_{\rm no})
=-\kappa A_{D_j(g)}(Q),\qquad
\kappa=\rot(P_{\rm loop})-\rot(P_{\rm no})\in\{-1,1\}.
\end{equation}
\status{proved (refereed: bench A, 2026-09-05T15:51:04Z; transcribed from RC mt:s4 (C022))}
\end{theorem}
\begin{proof}
Let $u=(\mu_{j-1},\mu_j)$, $v=(\mu_j,\mu_{j+1})$ and
$w=(\mu_{j-1},\mu_{j+1})$. These are physical edges; they are not
renumbered root slots. The cusp hypothesis places the deleted vertex
outside its neighbors' closed segment. As in the preceding theorem,
$\chi(\mu_{j-1},\mu_j,\mu_{j+1})=\tau_j$ by bilinearity.

First prove the signed consecutive identity directly, without assuming
the betweenness calculation from S3.  Write $A=\mu_{j-1}$, $M=\mu_j$,
$B=\mu_{j+1}$, and take line coordinates $A=0$, $B=1$, $M=t$,
where $t<0$ or $t>1$.  Name sides by $\chi(A,M,B)=-1,+1$ as in
S3.  The reverse-ordered far sign again gives $\delta=-1$ for
$P^+$ minus $P^-$.  For a nonincident root the two gaps are
singletons, regardless of their epsilon signs; hence $\mathcal U_L\mathcal U_R=1$.
The proper-span formula gives the signed deletion identity.

For the incident roots the following full-span table computes all four
cases from \eqref{afr:wall-epsilons}:
\begin{equation}\label{afr:s4-table}
\begin{array}{c|c|c|cc|c}
 \text{root}&\text{critical cut order}&t&
 \epsilon_L&\epsilon_R&
 \text{possibly nonzero summand}\\\hline
 u&M,B,A&t>1&1&1&\mathcal V_L\mathcal V_R=\mathcal B_R\\
 u&M,B,A&t<0&1&-1&\mathcal U_L\mathcal U_R=\mathcal B_R\\
 v&B,A,M&t>1&-1&1&\mathcal U_L\mathcal U_R=\mathcal B_L\\
 v&B,A,M&t<0&1&1&\mathcal V_L\mathcal V_R=\mathcal B_L
\end{array}
\end{equation}
For $u$, $L$ is singleton and $R$ is the deletion word closed by
$w=(A,B)$; the epsilon ratios are $(1-t)/(-t)$ and $(-1)/(-t)$.
For $v$, $R$ is singleton and $L$ is that deletion word; the
ratios are $(-1)/(t-1)$ and $t/(t-1)$.  These four ratio tests give
the table.  The other summand contains the nontrivial deletion
word's $E=0$.  Formula~\eqref{afr:wall-full} therefore gives
$A_g(P^+)-A_g(P^-)=-A_{D_j(g)}(Q)$ in each incident case also.

For the rotation sign, the principal local turn at $M$ tends to
$+\pi$ on the positive-chirotope side and to $-\pi$ on the
negative side, since its determinant is $\det(M-A,B-M)$ and
has sign $\chi(A,M,B)$.  Every other local turn is continuous at
this sole consecutive degeneration: its consecutive triple is distinct
from the critical triple because $n\geq4$, and is nonzero at the wall.
The edges are nonzero and the two incident directions at $M$ tend to
opposites. Thus the sum of all principal turns on the positive side
minus that on the negative side tends to $2\pi$. By
Lemma~\ref{lem:rot}, each punctured-side rotation is an integer constant
on its sufficiently small connected regular side. Their difference is
therefore its limit, namely $1$.  If the loop side is $P^+$, then
$\kappa=1$ and the just-proved signed formula is
\eqref{afr:S4}.  If it is $P^-$, then $\kappa=-1$ and
reversing the same formula gives \eqref{afr:S4}.
No root-independence or transport theorem is used.
\end{proof}

\begin{theorem}[vertex--edge law for $A_g$]\label{thm:A-S7}
At a simple vertex--edge wall at $(M;a)$, of bigon or sliding type,
let $s=\chi_{a,a+1,M}(P_-)$ and let $\lambda_1,\lambda_2$ be the
halves in Definition~\ref{def:deletion-halves}. For every physical root
$g$ with $H(g)=(h_1,h_2)$ as in Definition~\ref{def:induced-roots},
\begin{equation}\label{as7:law}
A_g(P_+)-A_g(P_-)=s A_{h_1}(\lambda_1)A_{h_2}(\lambda_2).
\end{equation}
\status{proved (refereed: bench A, 2026-09-05T15:51:04Z; transcribed from RC mt:s7 (C024))}
\end{theorem}
\begin{proof}
Write $A=\mu_a(0)$, $B=\mu_{a+1}(0)$, $X=\mu_M(0)$,
$d=(A,B)$, $d_1=(A,X)$ and $d_2=(X,B)$. The original cyclic order
of the three critical vertices is $A,B,X$. Both original arc words,
from $X$ to $A$ and from $B$ to $X$, have at least two leaves:
their closing halves are generic and have at least three vertices by
Lemma~\ref{lem:children}(ii). Each nontrivial closed word below has
only nonzero point-triple determinants. Its amplitude at the indicated
closing edge is therefore the interval value of Lemma~\ref{lem:farout}(iii).

In any root cut, the critical order is a cyclic permutation of $A,B,X$.
Such a permutation is even, whereas the reverse ordering in the far
sign is odd. The critical far sign is consequently
$-\chi(A,B,X)$. It equals $-s$ on $P_-$ and $s$ on $P_+$.
The quantity $\delta$ in \eqref{afr:wall-data} is therefore $s$.
Choose affine line coordinates $A=0$, $B=1$, $X=t$ with $0<t<1$.

For $g=d$, the boundary word begins at $B$, passes through $X$,
and ends at $A$. The critical span is full. The two ratios in
\eqref{afr:wall-epsilons} are $(t-1)/(-1)>0$ and $(-t)/(-1)>0$.
Both gaps are nontrivial, so both $\mathcal U$ factors are zero
$\mathcal E$ factors. The $\mathcal V$ factors are respectively
$A_{d_2}(\lambda_2)$ and $A_{d_1}(\lambda_1)$.
The full-span formula \eqref{afr:wall-full} gives the desired product.
Writing it in the prescribed half order uses commutativity of scalar
multiplication; it does not interchange the named halves.

If $g$ lies on the original arc from $X$ to $A$, the critical cut
order is $A,B,X$. The first gap is the singleton edge $d$, and the
second is the nontrivial arc from $B$ to $X$. Its epsilon ratios are
$1/t>0$ and $(t-1)/t<0$, so
$\mathcal U_L\mathcal U_R=A_{d_2}(\lambda_2)$.
The critical span is proper. Otherwise the physical root would be an
edge from $X$ to $A$, making that original arc a singleton and its
closing half a two-gon, contrary to the arity just established.
Contracting the span replaces the arc $A,B,\ldots,X$ by $d_1$.
The surviving polygon is exactly $\lambda_1$ with the same physical
root $g$. Formula~\eqref{afr:wall-proper} gives
$sA_g(\lambda_1)A_{d_2}(\lambda_2)$.

If $g$ lies on the original arc from $B$ to $X$, the critical cut
order is $X,A,B$. The first gap is the nontrivial arc from $X$ to
$A$, and the second is the singleton $d$. The ratios are
$-t/(1-t)<0$ and $1/(1-t)>0$. Thus
$\mathcal U_L\mathcal U_R=A_{d_1}(\lambda_1)$.
The span is proper, since a full span would make the original
$B$-to-$X$ arc a singleton. Contraction replaces $X,\ldots,A,B$
by $d_2$, leaving exactly $\lambda_2$ at the same physical root.
The proper-span formula gives $sA_{d_1}(\lambda_1)A_g(\lambda_2)$.

Every original root belongs to exactly one of these three cases.
The inherited-arc cases include the edges incident to the named
endpoints, so no additional root-seam case remains. No step restricts
the side of the contact line occupied by either neighbour of $X$.
Both bigon and sliding branches are therefore covered. The three
products are precisely the three rows of the stated half-root map.
\end{proof}

\begin{theorem}[triple law and silence for $A_g$]\label{thm:A-R3E}
\begin{enumerate}
\item[(i)] At a simple triple wall every composition weight, every
$b_{[i,j]}$ and $A_g$ are identical on the two sides, for every root $g$.
\item[(ii)] At a simple exterior-extension wall or a simple pure cut,
$A_g(P_+)=A_g(P_-)$ for every root $g$.
\end{enumerate}
\status{proved (refereed: bench A, 2026-09-05T15:51:04Z; transcribed from (ii) RC mt:extension, mt:pure; (i) as in SM1 (C024))}
\end{theorem}
\begin{proof}
For (i), every point-triple determinant has a nonzero central value,
so finitely many continuity conditions give a common interval on
which all signs read by Definition~\ref{def:gates} agree across the
wall. Every gate and composition weight agrees. Starting with the
one-leaf value and inducting on interval length in the open recursion
then gives equality of all $b$'s. Substituting them in the finite root
sum gives equality of $A_g$. Crossing-visit order is not an input to
any of these formulas.

For an exterior-extension wall, use $A,B,X,d$ as in the preceding
proof, but now $t<0$ or $t>1$. By the remoteness condition in
Definition~\ref{def:walls}(E), both full original arcs from $X$ to
$A$ and from $B$ to $X$ have at least two leaves. A singleton in
either would give a second adjacent pair in the critical triple and
would be a consecutive-triple wall instead. Closed nontrivial gap
and contracted words have (G1), since they omit at least one member
of the sole zero triple.

If $g=d$, the critical order is $B,X,A$, the span is full and both
gaps are nontrivial. Their epsilon ratios are $1-t$ and $t$.
Exactly one is negative. Hence $\mathcal U_L\mathcal U_R$ and
$\mathcal V_L\mathcal V_R$ each contain a nontrivial $\mathcal E$
factor, which is zero by Lemma~\ref{lem:farout}(iii).
The full-span formula gives zero jump.
If $g$ lies on the original $X$-to-$A$ arc, the cut order is $A,B,X$;
$L$ is singleton, $R$ is nontrivial, and
$\epsilon_R=\operatorname{sgn}((t-1)/t)=1$ for both allowed ranges
of $t$. Thus $\mathcal U_R=\mathcal E_R=0$. The span is proper,
since a full span would make the original $X$-to-$A$ arc a singleton.
Formula~\eqref{afr:wall-proper} gives zero jump.
For a root on the other arc, the cut order is $X,A,B$; $L$ is
nontrivial and $\epsilon_L=\operatorname{sgn}(-t/(1-t))=1$.
The span is proper for the analogous endpoint reason, and
$\mathcal U_L=\mathcal E_L=0$. This exhausts all exterior roots.
No assertion here applies to a consecutive cusp or flat wall.

Finally consider a pure cut. In the fixed root word write its critical
positions as $x<y<z$, with $N=n-1$. Pairwise nonadjacency gives
\begin{equation}\label{as7:pure-gaps}
y-x\geq2,\qquad z-y\geq2,\qquad N+1+x-z\geq2.
\end{equation}
The first two inequalities make both gaps nontrivial. The third shows
that the span cannot be full: full-span extremes would be the two
adjacent endpoints of the physical root. At least one epsilon is
positive, since the two nonzero line displacements from $x$ to $y$
and from $y$ to $z$ cannot both have sign opposite to their sum from
$x$ to $z$. Therefore the source product
$\mathcal U_L\mathcal U_R$ contains a zero $\mathcal E_L$ or
$\mathcal E_R$. Formula~\eqref{afr:wall-proper} again gives zero.
The argument allows $x=0$ or $z=N$ separately, so it includes roots
incident to a critical vertex and all six possible line orders.
The conclusion concerns the complete root output; no constancy of
the original open sums at a silent point-triple wall is asserted.
\end{proof}

\subsection{Continuation through silent configurations}

\begin{theorem}[continuation across silent walls]\label{thm:A-continuation}
For every root $g$ there is a unique function $\widetilde A_g$ on the weakly
generic locus $\mathcal W_n$ that is constant on visible chambers and agrees
with $A_g$ on $\mathcal U_n$.
\status{new (round~5: the inline openness and polygonal-connectedness arguments of the proof replaced by citations of Lemma~\ref{lem:weak-open}, F-25-147, and its density paragraph cites Lemma~\ref{lem:silentlocus} in place of a proof interior of Theorem~\ref{thm:relgp}, F-25-165; round~3: statement restricted to the continuation clause, F-25-68, F-25-71, F-25-89; the proof is otherwise the round-1 transcription of RC guard:unique, guard:covariance with the density and path argument as printed (C025), read in full by bench~A on SM2)}
\end{theorem}
\begin{proof}
First, the weakly generic locus $\mathcal W_n$ is open in the labelled
coordinate space, by Lemma~\ref{lem:weak-open}; in particular no new
crossing or triple concurrence appears in a small neighbourhood of a weakly
generic tuple, including at exterior collinearities.

The generic locus $\mathcal U_n$ is dense in the open set $\mathcal W_n$
(Lemma~\ref{lem:silentlocus}), so generic points occur in every open ball
contained in $\mathcal W_n$.

Every visible chamber, a connected component of the open set
$\mathcal W_n$, is polygonally connected, by Lemma~\ref{lem:weak-open}.

Take two generic points in one visible chamber $V$. Join them by a path
$\gamma$ in $V$. The compact image has a positive distance from the closed
complement of $\mathcal W_n$; choose a smaller positive error tolerance.
Since nonzero turns imply regularity, Theorem~\ref{thm:relgp} applies and
gives an endpoint-fixed path within that tolerance, hence still in
$\mathcal W_n$. Its possible actual wall events are F,V,T,E,C.
A flat event would have a zero turn, a vertex--edge event a vertex on a
nonincident closed segment, and a triple event a forbidden interior
concurrence. All three are excluded by membership in $\mathcal W_n$.
Only E and C remain. Along each generic segment $A_g$ is constant by
Proposition~\ref{prop:A-chamber}; at each of the finitely many remaining
events its jump is zero by Theorem~\ref{thm:A-R3E}(ii). Thus the two
generic endpoint values agree.

Density ensures that every visible chamber meets the generic locus.
Define $\widetilde A_g$ on that chamber to be the common generic value
just proved. This is a well-defined function, constant on visible chambers
and equal to $A_g$ on $\mathcal U_n$. Any other such function must have
the same value on each chamber because that chamber contains a generic
point. This proves uniqueness. The original finite tree formula on
$\mathcal U_n$ is the defining one; no evaluation of a gate at zero or
polynomial-form assertion at simultaneous silent zeros is used.
\end{proof}

\begin{remark}[the main text's amplitude on the silent locus]\label{rem:silent-amplitude}
On the silent locus $\mathcal S_n=\mathcal W_n\setminus\mathcal U_n$ the
tree sum of Definition~\ref{def:treesum} is not evaluated: its gates are
step functions of chirotope entries that vanish there, and no value
assigned to a gate at zero makes the literal sum independent of the root.
An exact six-gon with one exterior collinearity has literal root values
$(-1,-2t,-t,-t,-1,-1)$ when every gate at zero is given the value $t$,
while the continuation of Theorem~\ref{thm:A-continuation} takes the value
$-1$ at every root (phase ledger F-25-89: the Codex certificate, bench~A's
independent referee and the curator's replay). The main text's $A$ on its
generic locus, which is $\mathcal W_n$, is therefore \emph{defined} on
$\mathcal S_n$ as that continuation, in Section~\ref{sec:comparison}, once
root independence of the continuation is available; nothing in this
section assigns a value to a gate at zero.
\end{remark}


\begin{lemma}[multi-affine form]\label{lem:multiaffine}
Fix $n$ and a root index $g$ on labelled representatives. Replace
each cut factor in the finite tree sum by $(d-h)/2$ at an ordinary
vertex and $(d+h)/2$ at the root, interpreting each ordered
chirotope as the signed canonical variable $X_{ijk}$,
$1\leq i<j<k\leq n$. The resulting polynomial
$\mathsf A_g\in\mathbb Q[X_{ijk}:i<j<k]$ is affine in each variable
and satisfies $\mathsf A_g(\chi(P))=A_g(P)$ on (G1)-polygons.

Within a tree, each unordered vertex triple occurs in at most
one cut factor. At a binary vertex its near and far entries
coincide: the ordinary factor is zero, whereas the root factor
is that single entry. Distinct internal vertices can have the
same leaf interval only when they are the unary root and its
sole child; the unary root contributes an empty product.
\status{new (round~4: the binary-vertex case $d_1=h_1$ and the shared-interval case printed and the polynomial taken in the unrestricted ring $\mathbb Q[X_{ijk}]$; adopted from the Codex proposal sm4\_round4\_core\_v1, P-25-23, F115-POLY, re-read by the author; F-25-115; cleared by bench~A on SM5 (2026-09-06T00:32:58Z))}
\end{lemma}
\begin{proof}
The boundary word contains every vertex label exactly once.
Consequently unordered triples of boundary positions correspond
bijectively to unordered triples of vertex labels; alternating
the order changes only the sign of the canonical variable.

Consider a vertex with composition $r_0<\cdots<r_s$.
There are no cut factors when $s=1$. When $s=2$, the ordered
triples defining $d_1$ and $h_1$ are both
$(a_{r_2},a_{r_1},a_{r_0})$. Thus $d_1=h_1$ as formal
variables, and the ordinary and root factors are respectively
zero and $d_1$.

For $s\geq3$, the unordered near triple at cut $k$ is
$\{r_{k-1},r_k,r_{k+1}\}$ and the far triple is
$\{r_0,r_k,r_s\}$. Two near triples are equal only if their
smallest positions agree, which forces their cut indices to
agree. Two far triples are equal only if their interior
positions agree, again forcing the same cut. Equality of a
near triple and a far triple would require
$r_{k-1}=r_0$ and $r_{k+1}=r_s$ by comparison of their smallest
and largest positions. These equalities force $k=1$ and $s=2$,
contrary to $s\geq3$. Hence distinct cut factors at one vertex
use disjoint sets of canonical variables.

Now consider distinct vertices that both carry cut factors.
If neither is an ancestor of the other, their intervals lie
in different child intervals of their least common ancestor.
Those child intervals have disjoint interiors and at most one
common boundary position. Thus the two vertices cannot share
a triple of boundary positions.

If $v$ is an ancestor of $w$, the interval of $w$ lies in one
child interval of $v$. Each child interval of $v$ contains
exactly two cut positions of $v$, namely its endpoints.
Every near or far triple at $v$ contains three distinct cut
positions of $v$, so it cannot lie in that child interval
and cannot occur at $w$.

For completeness, a vertex with at least two children has
every child interval strictly smaller than its own.
All ordinary internal vertices have at least two children
by Lemma~\ref{lem:treesum-trees}. Therefore equality of the
intervals of distinct internal vertices is possible only
for the unary root and its sole child. The root then has no
cut factor, so this equality creates no repeated variable.

A tree with an ordinary binary vertex has zero formal weight.
In every other tree, the cut factors are linear and distinct
factors use disjoint sets of canonical variables. Every
monomial in their product therefore has degree at most one
in each variable. The sign $(-1)^{N_+(T)}$ does not change
this property, and neither does summing over the finite set
of trees. This proves multi-affinity in the polynomial ring,
without imposing any relations on its variables.

Finally, on (G1)-polygons Lemma~\ref{lem:gates-nonzero}
identifies every linear cut factor with the original gate.
Applying this identity term by term in
Lemma~\ref{lem:treesum-trees} proves
$\mathsf A_g(\chi(P))=A_g(P)$.
\end{proof}

\begin{lemma}[A nontrivial gap on a silent line]\label{pf:line-gap}
Let $a_0,\ldots,a_N$ be the boundary word of a weakly generic
polygon. Suppose $r_0<\cdots<r_k$, $k\geq2$, index collinear
vertices. Write their distinct line coordinates as $t_0,\ldots,t_k$
and put
\begin{equation}\label{pf:gap-directions}
 \epsilon_j=\sgn\frac{t_j-t_{j-1}}{t_k-t_0},\qquad 1\leq j\leq k.
\end{equation}
At least one gap has $\epsilon_j=1$ and $r_j-r_{j-1}\geq2$.
If $r_0=0$ and $r_k=N$, at least one gap also has
$\epsilon_j=-1$ and $r_j-r_{j-1}\geq2$.
\status{new (Codex candidate\_polyform\_dependent\_v1, C074; F-25-25 closed on SM2 by bench~A's referee of the two lemmas, R-25-21)}
\end{lemma}
\begin{proof}
The vertices of a weakly generic polygon are distinct. Coincident
adjacent vertices would give a zero edge; coincident nonadjacent
vertices would put a vertex on a nonincident closed edge.
An original edge whose endpoints are among the selected line points
cannot contain any other selected point in its interior. Also two
successive gaps of one leaf cannot both join selected line points:
they would be consecutive collinear original edges, with zero turn.
These are consequences of Definition~\ref{def:weak}.

Choose the line coordinate so $t_0=0$ and $t_k=1$. Suppose every
increasing gap had one leaf. Let $t_\ell$ be the largest selected
coordinate; it is not $t_0$, so $\ell\geq1$. Its incoming gap is
increasing, hence is an original edge. Since that edge contains no
other selected point, $t_{\ell-1}$ is the second largest selected
coordinate. We cannot have $\ell=1$: then this second largest
coordinate would be $t_0=0$, whereas the distinct final point has
coordinate $1$ below the maximum. The maximum cannot itself be
that final point when $\ell=1$, since $k\geq2$.
Thus $\ell\geq2$. The preceding coordinate $t_{\ell-2}$ is smaller
than $t_{\ell-1}$, because only $t_\ell$ is larger and the selected
points are distinct. That preceding gap is therefore also an
increasing original edge. The two successive one-leaf gaps give
the forbidden zero turn. This proves the first assertion.

For the second assertion close the selected list by the physical
root edge from $a_N$ to $a_0$. Its line direction is negative.
If every negative gap in the open list had one leaf, every negative
step in this cyclic list would be an original edge. Start at its
largest coordinate. The following point is lower, and the edge to
it cannot contain another selected point, so it is the second largest
coordinate. There are at least three selected points. The next point
in cyclic order is therefore still lower: it is neither of those two
points. The next step is another negative original edge, again
giving two successive collinear original edges and a zero turn.
This contradiction proves that one negative open gap is nontrivial.
\end{proof}

\begin{lemma}[Formal cancellation of all silent entries]
\label{pf:formal-cancellation}
For the boundary word of any $P\in\mathcal W_n$, let $H^0$ be its
geometric far array, with its zero entries equal to zero. Form $H$
by replacing every zero entry of $H^0$ by its own indeterminate,
keeping all nonzero entries fixed. Over the polynomial ring in those
indeterminates, put $c=F_H^{-1}(E)$ and $c^0=F_{H^0}^{-1}(E)$.
Then each coordinate $c_{ij}$ equals the corresponding constant $c^0_{ij}$, and
\begin{equation}\label{pf:formal-output}
 F_{-H}(c)_{0N}=F_{-H^0}(c^0)_{0N}.
\end{equation}
No independence of the zero-determinant differentials is assumed.
\status{new (Codex candidate\_polyform\_dependent\_v1, C074; F-25-25 closed on SM2 by bench~A's referee of the two lemmas, R-25-21)}
\end{lemma}
\begin{proof}
The polynomial inverses exist by Lemma~\ref{lem:farout}(i).
We first establish a coefficient identity with the child coordinates
fixed to $c^0$. Fix $I=[i,j]$ and let
$Z_I=\{r:i<r<j,\ H^0(i,r,j)=0\}$. Only the entries indexed by
$r\in Z_I$ vary in this top-coordinate formula. Write those entries
as $u_r$. For a nonempty subset
$S=\{r_1<\cdots<r_{k-1}\}\subseteq Z_I$, set $r_0=i$, $r_k=j$
and $J_\ell=[r_{\ell-1},r_\ell]$. All these endpoints lie on the
line through $a_i,a_j$; define $\epsilon_\ell$ as
in~\eqref{pf:gap-directions}. For $\eta\in\{1,-1\}$ the coefficient
of $\prod_{r\in S}u_r$ in the top formula is
\begin{equation}\label{pf:coefficient-factor}
 \left[\prod_{r\in S}u_r\right]F_{\eta H}(c^0)_{ij}
 =\left(-\frac{\eta}{2}\right)^{|S|}
   \prod_{\ell=1}^{k}
      F_{\eta\epsilon_\ell H^0}(c^0)_{J_\ell}.
\end{equation}
Here the bracket denotes coefficient extraction only in the top
variables $u_r$; every child coordinate is the fixed number $c^0$.

To prove the identity, each top cut occurs at most once in a
composition. A contribution to this coefficient must include every
cut in $S$ and exclude all other cuts in $Z_I$. Split its ordered
cut list at the members of $S$. This is a bijection with independent
compositions of the gaps $J_\ell$, excluding their other collinear
cuts; concatenation reverses it. Each selected cut supplies
$-\eta/2$. At a further cut $r$ in gap $J_\ell$, collinearity of
the gap endpoints gives
\begin{equation}\label{pf:gap-sign-identity}
 H^0(i,r,j)=\epsilon_\ell
 H^0(r_{\ell-1},r,r_\ell).
\end{equation}
Indeed write the endpoints on the common line as $p+t\omega$.
In the determinant defining $H^0$, changing the endpoint difference
from $(t_k-t_0)\omega$ to
$(t_\ell-t_{\ell-1})\omega$ multiplies its sign by
$\epsilon_\ell$; translating the line endpoint changes no
determinant with $\omega$. This calculation also holds when $a_r$
lies on that line: both signs are then zero. Such a cut has zero
weight in the gap transform, so including it there changes no sum.
The remaining cuts and child factors are exactly those of the gap
transforms in~\eqref{pf:coefficient-factor}. Finite distributivity
now proves that identity, including gaps with only one leaf.

For $\eta=1$, Lemma~\ref{pf:line-gap} supplies a nontrivial gap
with $\epsilon_\ell=1$. Its factor is
\begin{equation}\label{pf:zero-gap-factor}
 F_{H^0}(c^0)_{J_\ell}=E_{J_\ell}=0.
\end{equation}
The first equality is the defining inverse equation, and the second
uses the gap's length of at least two. Thus every nonempty top
coefficient vanishes. The constant coefficient is the value with
all $u_r=0$, namely $F_{H^0}(c^0)_{ij}=E_{ij}$.

Induct on interval length to prove $c=c^0$ for the entire formal
array $H$. Leaf coordinates are one in both arrays. If all shorter
coordinates agree, the equation for $[i,j]$ has its unknown unary
term $c_{ij}$ with coefficient one; every nonunary child is already
$c^0$. The preceding coefficient calculation says that substituting
also $c^0_{ij}$ makes this equation equal $E_{ij}$ for every value
of the top variables. Uniqueness of its unary solution therefore
gives $c_{ij}=c^0_{ij}$. This proves the induction. In particular it
does not require the formal sign assignments to be realizable.

Finally apply~\eqref{pf:coefficient-factor} to $I=[0,N]$ and
$\eta=-1$. For every nonempty $S$, the second assertion of
Lemma~\ref{pf:line-gap} supplies a nontrivial gap with
$\epsilon_\ell=-1$. Its factor is again $F_{H^0}(c^0)_{J_\ell}=0$.
All nonconstant top coefficients vanish, and the constant term is
$F_{-H^0}(c^0)_{0N}$. Substituting $c=c^0$ proves
\eqref{pf:formal-output}. The physical root edge is used here;
an arbitrary open interval would not justify this last negative-gap
step.
\end{proof}

\begin{corollary}[The continuation is the polynomial form]
\label{cor:polyform}
On all of $\mathcal W_n$, for every root $g$,
\begin{equation}\label{pf:polynomial-value}
 \widetilde A_g(P)=\mathsf A_g(\chi(P)).
\end{equation}
After the nonzero chirotope entries are fixed at their values at $P$,
the polynomial is constant in all remaining entries. Thus its
nonconstant terms in these silent variables cancel, and no value
of $\Theta(0)$ has to be chosen.
\status{new (round~6: the specialization stated at the cut-factor granularity of Definition~\ref{def:gates} --- each gate variable is sent to the product over the interior cuts of its composition of the linear forms of Lemma~\ref{lem:multiaffine}, F-25-166; round~5: the recursion-built formal output identified with the tree-sum polynomial $\mathsf A_g$ in $\mathbb Q[X_{ijk}]$ through the formal clause of Lemma~\ref{lem:treesum-trees}, F-25-150; round~4: the openness and density citations of the ball step repaired, Lemma~\ref{lem:weak-open} and Lemma~\ref{lem:silentlocus}, F-25-116; Codex candidate\_polyform\_dependent\_v1, C074; unrestricted argument; F-25-25 closed on SM2, R-25-21; assembly: Lemmas~\ref{pf:line-gap} and~\ref{pf:formal-cancellation} refereed (R-25-21), Lemma~\ref{lem:treesum-trees} new (round~5, its formal clause), Lemma~\ref{lem:multiaffine} and Lemma~\ref{lem:weak-open} new (round~4), Theorem~\ref{thm:A-continuation} new (read in full by bench~A on SM2), Lemmas~\ref{lem:silentlocus} and~\ref{lem:farout} proved; the evaluation of the continuation is in Section~\ref{sec:comparison})}
\end{corollary}
\begin{proof}
Replace each cut factor of every gate in the finite tree recursion by
the linear form $(d_k-h_k)/2$ or $(d_k+h_k)/2$ of
Lemma~\ref{lem:multiaffine}, and regard the entries as formal
variables. The same recursion is
$\Phi_{D,H}(b)=E$, with root output $\Phi_{D,-H}(b)_{0N}$.
By the formal clause of Lemma~\ref{lem:treesum-trees}, the recursion and the
tree sum agree as polynomials in the gate variables; sending each gate
variable $V^\pm(\pi)$, for a composition $\pi=(r_0<\cdots<r_s)$, to the
product over the interior cuts $r_k$, $1\leq k\leq s-1$, of the linear forms
of Lemma~\ref{lem:multiaffine} --- $(d_k-h_k)/2$ for $V^+$ at an ordinary
vertex and $(d_k+h_k)/2$ for $V^-$ at the root, the empty product being $1$
--- defines a ring homomorphism from $\ZZ[V^\pm(\pi)]$ to
$\ZZ[\tfrac12][X_{ijk}]\subset\mathbb Q[X_{ijk}]$ (Definition~\ref{def:gates}
makes each gate that product of its cut factors), so this root output is the
polynomial $\mathsf A_g$ of that lemma, as an identity in
$\mathbb Q[X_{ijk}]$ and not only at realizable chirotopes.
The ring identity of Lemma~\ref{lem:farout}(i), followed by its
two-sided inverse, makes this formal output
$F_{-H}(F_H^{-1}(E))_{0N}$. This is a polynomial identity, so it
also applies when the geometric signs have zeros; it does not
evaluate the original step-function gates at zero.
The geometric near and far entries are the reverse-ordered
chirotope entries of the root word. This changes only fixed signs
and names of the canonical variables. Lemma~\ref{pf:formal-cancellation}
therefore makes $\mathsf A_g$ constant in every entry that vanishes
at $P$, with the other entries fixed.

Fix the labelled representative $P$ and root $g$. By
Lemma~\ref{lem:weak-open}, choose a ball about $P$ contained in its
visible chamber, small enough that every initially nonzero
area determinant keeps its sign. By Lemma~\ref{lem:silentlocus},
this ball contains a generic tuple $Q$. On $Q$ the polynomial equals $A_g(Q)$
by Lemma~\ref{lem:multiaffine}, and this equals
$\widetilde A_g(P)$ by Theorem~\ref{thm:A-continuation}.
Formal constancy gives the same value after all silent entries are
set to zero, which is $\mathsf A_g(\chi(P))$. This proves the
identity and the cancellation assertion throughout the weak locus.
\end{proof}


\subsection{The soft theorem}

\begin{definition}[soft insertion]\label{def:soft}
Let $P$ satisfy (G1), fix a vertex index $j$ and a vector $q\neq0$. The
\emph{soft insertion} of $q$ at $j$ is the family of $(n+1)$-gons
\[
P^{(j,q)}_\varepsilon=(\mu_1,\ldots,\mu_j,\ \mu_j+\varepsilon q,\ \mu_{j+1},\ldots,\mu_n),
\qquad\varepsilon>0,
\]
with the new vertex $\mu_*=\mu_j+\varepsilon q$ inserted immediately after
$\mu_j$; its edges are $\lt_{j-1}$, the \emph{soft edge} $\varepsilon q$, and
the \emph{return edge} $\lt_j-\varepsilon q$, the others unchanged. The
\emph{attachment signs} are
\[
\chi_-=\chi_{*,j,j-1}\bigl(P_\varepsilon\bigr)=-\sgn\det(\lt_{j-1},q),\qquad
\chi_+=\chi_{*,j,j+1}\bigl(P_\varepsilon\bigr)=-\sgn\det(q,\lt_j),
\]
the two equalities holding for every $\varepsilon>0$ (they are the turns at
the two ends of the soft edge, with the sign reversed). The vector $q$ is
\emph{admissible} if $\det(\lt_{j-1},q)\neq0$, $\det(q,\lt_j)\neq0$, and
$\det(q,\mu_k-\mu_j)\neq0$ for all $k\neq j$.
\end{definition}

\begin{lemma}[the soft family is generic]\label{lem:soft-generic}
Let $P$ be generic, $j$ a vertex, and $q$ admissible in
Definition~\ref{def:soft}. Put $M=\mu_j$, $A=\mu_{j-1}$,
$B=\mu_{j+1}$, $u=M-A$, $v=B-M$, and $\tau=\tau_j(P)$.
Write $M_\varepsilon=M+\varepsilon q$ and
$D_\varepsilon=v-\varepsilon q$. There is $\varepsilon_0>0$
such that, for every $0<\varepsilon<\varepsilon_0$:
\begin{enumerate}
\item[(i)] $P_\varepsilon=P^{(j,q)}_\varepsilon$ is generic,
and all these polygons lie in one chamber. The soft edge meets
other edges only at its incident endpoints.
\item[(ii)] Every original vertex other than $M$ retains its
original turn sign. The turns at $M$ and $M_\varepsilon$ are
$-\chi_-$ and $-\chi_+$, respectively. The unchanged edge
 directions are fixed, $D_\varepsilon\to v$, and
$\sgn\det(u,D_\varepsilon)=\tau$, while
$\sgn\det(D_\varepsilon,u)=-\tau$.
\item[(iii)] Identify each unchanged edge with its parent edge
and the return edge with the parent edge $E_j$. Every parent
crossing persists. Its point and its parameters on the corresponding
edges converge to the parent values. Inherited visits retain their
order along each directed edge, and the determinant sign of the
ordered edge directions at each inherited crossing is unchanged.
\item[(iv)] If $\chi_-=\chi_+=-\tau$ (same-sign sector), or
$\chi_-\ne\chi_+$ (mixed sector), these are all the crossings
and the Gauss word is the parent's Gauss word. If
$\chi_-=\chi_+=\tau$ (loop sector), there is exactly one
additional crossing $y$, between the return edge and $E_{j-1}$.
It tends to $M$. Its incoming-edge visit $a$ and return-edge
visit $b$ bound the oriented traversal arc from $a$ through
$M,M_\varepsilon$ to $b$, containing no other crossing visit.
Thus the two newborn visits are cyclically adjacent; deleting
 them gives the parent's Gauss word.
\end{enumerate}
\status{new (round~4: the geometric facts consumed by Theorem~\ref{thm:C-soft} exported as clauses (i)--(iv) --- turn signs, soft-edge contacts, inherited crossings with their parameters, order and direction-determinant signs, the unique newborn crossing and its short arc; adopted from the Codex proposal sm4\_round4\_exports\_v1, P-25-24, F120-GEOMETRY, re-read by the author; F-25-120; cleared by bench~A on SM5 (2026-09-06T00:32:58Z))}
\end{lemma}
\begin{proof}
Retain the original names of vertices and edges, replacing only
$E_j$ by the soft and return edges. The positive parameter interval
may be shortened finitely many times in the argument.

First establish (G1). Original vertex triples are unchanged.
For distinct original indices $k,l\ne j$, the area determinant for
$(M_\varepsilon,\mu_k,\mu_l)$ tends to the nonzero determinant
for $(M,\mu_k,\mu_l)$. For $k\ne j$, the determinant for
$(M_\varepsilon,M,\mu_k)$ is
$-\varepsilon\det(q,\mu_k-M)$, nonzero by admissibility.
These exhaust the new triples. Finiteness gives one interval
on which (G1) holds and all signs are constant.

The soft edge lies within $\varepsilon|q|$ of $M$.
Every parent segment not incident to $M$ has positive distance
from $M$, by compactness and (G1) of $P$. Hence the soft edge
misses all such segments for small $\varepsilon$. It meets the
incoming edge only at $M$, since $\det(u,q)\ne0$, and the return
edge only at $M_\varepsilon$, since
$\det(q,D_\varepsilon)=\det(q,v)\ne0$. The return edge and the
unchanged edge following $B$ meet only at $B$: their directions
tend to the independent directions at the parent corner $B$.
All pairs of unchanged edges retain their parent intersections.

For each parent edge remote to $E_j$, compare it with the return
edge. A disjoint parent pair has positive compact distance;
the return segment converges uniformly at each affine parameter
to $E_j$, so that distance remains positive. No independence
of the directions is required for a disjoint pair. For a crossing
parent pair, the directions have nonzero determinant and both
intersection parameters lie in $(0,1)$. The intersection equations
form an invertible two-by-two linear system. Its coefficients
converge to the parent coefficients and its determinant remains
nonzero, so its inverse varies continuously. The two parameters
converge to the parent values and stay in $(0,1)$; the point
converges as well, and the ordered direction determinant keeps
its sign.

These are all return-edge pairs except the incoming edge
$E_{j-1}$. Put $T=\det(u,v)$, of sign $\tau$. The return
edge meets $E_{j-1}$ precisely when each supporting line
strictly separates the endpoints of the other segment; the
directions are independent since $\det(u,D_\varepsilon)\to T\ne0$.
For the incoming line the signed heights are
\begin{equation}\label{eq:soft-first-straddle}
\det(u,M_\varepsilon-A)=\varepsilon\det(u,q),\qquad
\det(u,B-A)=T.
\end{equation}
They are opposite exactly when $\chi_-=\tau$. For the return
line based at $M_\varepsilon$, the height of $M$ is
\begin{equation}\label{eq:soft-second-straddle}
\det(D_\varepsilon,M-M_\varepsilon)
=\varepsilon\det(q,v).
\end{equation}
The height of $A$ tends to $\det(v,-u)=T\ne0$, so has sign
$\tau$ throughout a sufficiently small interval. This pair of
heights is opposite exactly when $\chi_+=\tau$. Both tests
are required: their conjunction is exactly the loop sector.
There is no intersection of this pair in the other two sectors.
In the loop sector its supporting-line equations remain
invertible as $\varepsilon\to0$. At zero the unique intersection
is $M$, with incoming parameter one and return parameter zero.
Thus $y\to M$, and these two parameters tend to one and zero.

All parent crossing points are away from $M$, since (G1)
excludes a vertex on a nonincident edge, and they are pairwise
distinct by (G2). The finitely many inherited points therefore
remain pairwise distinct and away from $M$. The newborn, when
present, stays distinct from them after another shrink.
Every possible crossing has been enumerated. A concurrence of
three distinct edge interiors would involve pairwise remote
edges, since (G1) makes distinct adjacent interiors disjoint.
It would identify the points of two distinct unordered crossing
pairs, which was excluded. Hence (G2) holds. The connected
parameter interval now maps continuously into the labelled
generic locus, so its image lies in one labelled chamber and
hence one polygon chamber. This proves (i).

The turn determinants at $M$ and $M_\varepsilon$ are
$\varepsilon\det(u,q)$ and
$\varepsilon\det(q,D_\varepsilon)=\varepsilon\det(q,v)$;
Definition~\ref{def:soft} gives their signs $-\chi_-$ and
$-\chi_+$. At $B$, the incoming direction tends to $v$ and
the outgoing one is unchanged, so the old nonzero turn sign
persists. Every other old vertex except $M$ has unchanged
incident directions. Finally $\det(u,D_\varepsilon)\to T$
gives sign $\tau$, and antisymmetry gives $-\tau$ for the
opposite ordered pair. This proves (ii).

Distinct parent crossing parameters on each edge have positive
pairwise separations and positive distances from zero and one.
Their convergence preserves their order on the corresponding
edges, including every edge whose crossing points move.
This proves (iii). Without a newborn, concatenating these
edgewise visit lists recovers the parent's cyclic Gauss word.
With a newborn, its incoming visit is after all inherited visits
on the incoming edge, and its return visit is before all inherited
visits on the return edge. The intervening soft edge has no crossing.
Thus the arc from $a$ through $M,M_\varepsilon$ to $b$ has no
other crossing visit. Deleting $a,b$ leaves exactly the inherited
lists, proving (iv).
\end{proof}

\begin{theorem}[soft theorem for $A_g$]\label{thm:A-soft}
Let $P$ be generic, $q$ admissible, and $\varepsilon_0$ as in
Lemma~\ref{lem:soft-generic}. For every root $g$ of $P_\varepsilon$ other than
the soft edge, with $g'$ the corresponding root of $P$ (the return edge
corresponds to $E_j$), there is $\varepsilon_1\in(0,\varepsilon_0]$ with
\begin{equation}\label{aspr:soft-law}
A_g(P_\varepsilon)=\frac{\chi_-+\chi_+}{2}A_{g'}(P)
\qquad(0<\varepsilon<\varepsilon_1).
\end{equation}
Equivalently, in the same-sign, mixed and loop sectors the multipliers are
$-\tau_j(P)$, $0$ and $\tau_j(P)$, respectively.
\status{proved (refereed: bench A, 2026-09-05T15:51:04Z; transcribed from RC mt:pruning, mt:soft-duplication (C017))}
\end{theorem}
\begin{proof}
We prove the finite duplication identity used below, including its two
boundary cases. This proof uses the definitions of the boundary word,
gates, open sums and soft insertion. It does not use a wall law or root
independence. All formal arrays in the next two paragraphs are over
$\mathbb Q$; their near entries need not be geometric signs.

\emph{Finite transforms and freedom to choose a near array.}
For a linearly ordered list with boundary positions $0,\ldots,L$, define
\begin{equation}\label{aspr:Phi}
\Phi_{D,H}(X)_{ij}=
\sum_{\pi=(i=r_0<\cdots<r_s=j)}
\prod_{k=1}^{s-1}
\frac{D(r_{k-1},r_k,r_{k+1})-H(i,r_k,j)}2
\prod_{k=1}^sX_{r_{k-1},r_k}.
\end{equation}
Put $F_H=\Phi_{0,H}$, $G_D=\Phi_{D,0}$, and let $E_{ij}$ be $1$ when
$j=i+1$ and $0$ otherwise. The one-part composition contributes $X_{ij}$;
every other term uses only shorter intervals. Thus one solves
$\Phi_{D,H}(X)=Y$ uniquely in increasing interval length, by subtracting
the already determined nonunary terms from $Y_{ij}$. This gives a
polynomial inverse on the finite array space. The constructed solution
map is a right inverse; uniqueness applied to the target
$\Phi_{D,H}(X)$ makes it a left inverse too.

The following factorization holds for arbitrary near and far arrays:
\begin{equation}\label{aspr:factorization}
\Phi_{D,H}=F_H\circ G_D.
\end{equation}
Indeed, expand the right side using an outer composition of $[i,j]$ and
an inner composition of each outer part. Their ordered union is a refined
composition $\rho$, together with a subset of its interior cuts marked as
outer cuts. Conversely, that subset recovers the outer composition, and
restricting $\rho$ to each part recovers the inner compositions. These
operations are inverse, including empty or full outer-cut subsets and
one-part inner or outer compositions. At an outer cut $r_k$ the factor is
$-H(i,r_k,j)/2$. At a nonouter cut its immediate neighbors in $\rho$
are exactly its neighbors inside its inner part, so the factor is
$D(r_{k-1},r_k,r_{k+1})/2$. Each interior cut supplies one factor, and the
child-interval product is the refined product. Summing independently over
the outer-cut subset gives the product of their sums, which is
\eqref{aspr:Phi}.

Write $b=\Phi_{D,H}^{-1}(E)$ and $c=F_H^{-1}(E)$. Applying
\eqref{aspr:factorization} to the equation for $b$ gives
$F_H(G_D(b))=E$; uniqueness gives $G_D(b)=c$. Applying the same
factorization with $-H$ then gives
\begin{equation}\label{aspr:near-free}
\Phi_{D,-H}(b)=F_{-H}(c).
\end{equation}
Thus the complete root transform is independent of $D$ with $H$ fixed.
For geometric arrays
$D(x,y,z)=H(x,y,z)=\chi(a_z,a_y,a_x)$, the elementary gate identities
$d\Theta(-hd)=(d-h)/2$ and $d\Theta(hd)=(d+h)/2$ show that
$\Phi_{D,H}(b)=E$ is exactly the open-sum recursion, and its full
$\Phi_{D,-H}(b)$ coordinate is $A_g$, including the one-part root term.
The freedom in \eqref{aspr:near-free} concerns the complete transform;
it asserts no equality of individual auxiliary and geometric open sums.

\emph{The formal duplication problem.}
Take a core ordered list $z_0,\ldots,z_M$, $M\geq2$, and distinguish
$A=z_s$, where $0\leq s\leq M$. Insert a new occurrence $B$ immediately
after $A$. The collapse $r$ sends $B$ to $A$ and fixes the other
occurrences. Let $H^0$ be any core far array of signs, and let
$t(X)\in\{-1,1\}$ be arbitrary for each core occurrence $X\ne A$.
Assume the parent far array $H^1$ has the following entries, where order
relations refer to positions in the linear lists:
\begin{equation}\label{aspr:far-data}
\begin{array}{ll}
H^1\text{ on a triple not containing both }A,B
   &\text{is the pullback of }H^0,\\
H^1(X,A,B)=t(X)&(X<A),\\
H^1(A,B,Y)=t(Y)&(B<Y).
\end{array}
\end{equation}
Only the neighbors in the next display are read cyclically in the core:
\begin{equation}\label{aspr:neighbor-data}
\eta_-=t(z_{s-1\bmod(M+1)}),\qquad
\eta_+=t(z_{s+1\bmod(M+1)}),\qquad
k=\frac{\eta_-+\eta_+}{2}.
\end{equation}
We will prove that the full parent root output is $k$ times the core one.

Choose auxiliary near arrays as follows. On core triples with middle
occurrence $A$, set $D^0(X,A,Y)=t(X)$; on other core triples set $D^0=0$.
On parent triples not containing both $A,B$, pull back $D^0$, and set
\begin{equation}\label{aspr:near-data}
D^1(X,A,B)=t(X)\ (X<A),\qquad
D^1(A,B,Y)=t(Y)\ (B<Y).
\end{equation}
These are formal arrays, not a claimed chirotope. Keep the far arrays
fixed. Let $b^0=\Phi_{D^0,H^0}^{-1}(E^0)$, and propose the following
parent ordinary array:
\begin{equation}\label{aspr:ordinary-table}
\begin{array}{c|c}
\text{parent interval }I&b^1_I\\\hline
I\text{ not containing both }A,B&b^0_{r(I)}\\
{[A,B]}&1\\
{[A,Y]},\ B<Y&\dfrac{\eta_+-t(Y)}2\,b^0_{[A,Y]}\\
{[X,B]},\ X<A&\dfrac{\eta_--t(X)}2\,b^0_{[X,A]}\\
{[X,Y]},\ X<A<B<Y&k\,b^0_{[X,Y]}
\end{array}
\end{equation}
Consecutiveness of $A,B$ makes these cases exhaustive. Apart from the
explicit singleton $[A,B]$, the collapsed endpoints are distinct. We
verify $\Phi_{D^1,H^1}(b^1)=E^1$ on every interval before identifying
$b^1$ with the actual auxiliary ordinary solution.

\emph{Intervals without the duplicate and the singleton.}
On an interval not containing both occurrences, collapse is an
order-preserving bijection on vertices, compositions, gates and child
intervals. Thus its transform is the corresponding core $E^0$ entry,
which equals its parent $E^1$ entry. On $[A,B]$ the sole composition has
one part and value $1$, as required.

\emph{Intervals starting at $A$.}
Fix $[A,Y]$ with $B<Y$. For a core composition of $[A,Y]$, write $X$
for its first boundary after $A$. There are exactly two parent
presentations. Without a cut at $B$, its first child contains $A,B$ and
has multiplier $(\eta_+-t(X))/2$. With a cut at $B$, the first child is
the singleton $[A,B]$, and the rest is the core composition with its
first endpoint renamed $B$. Its new gate at $B$ has near entry $t(X)$
and far entry $t(Y)$. Removing the cut at $B$, when present, recovers
the core composition; its presence distinguishes the two inverse
expansions. All other children agree, the neighboring near gate pulls
$B$ back to $A$, and all other far entries are the stipulated pullbacks.
The sum of the two local multipliers is, for an ordinary top vertex,
\begin{equation}\label{aspr:first-ordinary}
\frac{\eta_+-t(X)}2+\frac{t(X)-t(Y)}2
=\frac{\eta_+-t(Y)}2,
\end{equation}
and for a root top vertex it is
\begin{equation}\label{aspr:first-root}
\frac{\eta_+-t(X)}2+\frac{t(X)+t(Y)}2
=\frac{\eta_++t(Y)}2.
\end{equation}
In each case $t(X)$ cancels, leaving a composition-independent multiplier
of the corresponding core transform. The ordinary transform vanishes
when the core has at least two leaves, since its $E^0$ entry is zero.
When the core has one leaf, $Y$ is the actual successor of $A$, so
$t(Y)=\eta_+$ and the ordinary multiplier itself vanishes. Thus every
such parent ordinary coordinate has its required zero value, including
the two-leaf parent interval.

\emph{Intervals ending at $B$.}
Fix $[X,B]$ with $X<A$. In a core composition of $[X,A]$, let $Y$
be the last boundary before $A$. Without a cut at $A$ the last child
contains the soft edge and has multiplier $(\eta_--t(Y))/2$.
With that cut it ends in the singleton $[A,B]$; the new gate has near
entry $t(Y)$ and far entry $t(X)$. Remove the cut $A$ when present
and rename the endpoint $B$ as $A$ to recover the core list. The cut's
presence distinguishes the two expansions. The neighboring near gate
pulls $B$ back to $A$ and all other factors agree. Consequently the
ordinary and root multipliers are, respectively,
\begin{equation}\label{aspr:last-ordinary}
\frac{\eta_--t(Y)}2+\frac{t(Y)-t(X)}2
=\frac{\eta_--t(X)}2,
\end{equation}
\begin{equation}\label{aspr:last-root}
\frac{\eta_--t(Y)}2+\frac{t(Y)+t(X)}2
=\frac{\eta_-+t(X)}2.
\end{equation}
For a core interval with at least two leaves its ordinary transform is
zero. For one leaf, $X$ is the actual predecessor of $A$, so
$t(X)=\eta_-$ and the multiplier is zero. This checks the other
two-leaf parent boundary as well.

\emph{Intervals strictly spanning $A,B$.}
Fix $[X,Y]$ with $X<A<B<Y$. Collapse every parent cut list, retaining
only one copy when both $A,B$ are cuts. If the resulting core
composition does not cut at $A$, it has exactly one child spanning $A$.
It has exactly one parent presentation, with neither $A$ nor $B$ a cut;
that child strictly spans the duplicate and has multiplier $k$.
All top gates pull back unchanged. This includes the one-part
composition: the proposed leading $b^1$ coordinate is being checked,
not assumed to satisfy a recurrence.

If the core composition cuts at $A$, let $L,R$ be its neighboring cuts
and put $a=t(L)$, $b=t(R)$ and $h=H^0(X,A,Y)$. Its core near entry at
$A$ is $a$. There are precisely three parent presentations:
\begin{equation}\label{aspr:three-presentations}
\begin{array}{c|c}
\text{cuts among }A,B&\text{exceptional child multiplier}\\\hline
B\text{ only}&(\eta_--a)/2\\
A\text{ only}&(\eta_+-b)/2\\
A\text{ and }B&1\quad([A,B]\text{ is a singleton})
\end{array}
\end{equation}
Collapse recovers the core composition, while its subset of cuts among
$A,B$ recovers the row; expanding according to each row is the inverse.
The fourth subset, neither cut, gives the preceding no-core-cut case
and cannot occur here. The neighboring gate at $L$ has right neighbor
$A$ or $B$, both pulling back to $A$; the neighboring gate at $R$
has the analogous left neighbor. Their near entries are unchanged.
The far entries at $A$ and $B$ both pull back to $h$, and all exterior
gates and remaining child factors agree with the core.

For an ordinary top, either one-cut row has the core gate $(a-h)/2$;
the two-cut row replaces it by $(a-h)(b-h)/4$. The two one-cut child
multipliers sum to $k-(a+b)/2$. After subtracting $k$ times the core
gate, the residual is therefore
\begin{equation}\label{aspr:ordinary-residual}
R_{\rm o}=-\frac{a+b}{2}\frac{a-h}{2}
               +\frac{(a-h)(b-h)}4.
\end{equation}
Factoring $(a-h)/4$ leaves $-(a+b)+(b-h)=-a-h$, so
\begin{equation}\label{aspr:ordinary-factor}
R_{\rm o}=\frac{(a-h)(-a-h)}4=\frac{h^2-a^2}{4}=0.
\end{equation}
The last equality uses exactly $a,h\in\{-1,1\}$.
For a root top the corresponding residual is
\begin{equation}\label{aspr:root-residual}
R_{\rm r}=-\frac{a+b}{2}\frac{a+h}{2}
               +\frac{(a+h)(b+h)}4.
\end{equation}
Factoring $(a+h)/4$ leaves $-(a+b)+(b+h)=-a+h$, so
\begin{equation}\label{aspr:root-factor}
R_{\rm r}=\frac{(a+h)(-a+h)}4=\frac{h^2-a^2}{4}=0.
\end{equation}
Thus both transforms equal $k$ times their respective core transforms
on every strictly spanning interval. Such a core interval has at least
two leaves, so its ordinary $E^0$ entry is zero. No relation among
different remote values of $t$ was used.

All interval cases have now verified
$\Phi_{D^1,H^1}(b^1)=E^1$. Finite triangular uniqueness proved above
identifies \eqref{aspr:ordinary-table} with the actual auxiliary
ordinary solution. The root transforms already checked therefore
compute the actual auxiliary root outputs. We did not assume that any
auxiliary binary open sum vanishes.

For the full interval, if $0<s<M$ the strictly spanning multiplier is
$k$. If $s=0$, the final endpoint $Y=z_M$ is the cyclic predecessor of
$A$, so $t(Y)=\eta_-$ and \eqref{aspr:first-root} gives $k$. If $s=M$,
the initial endpoint $X=z_0$ is the cyclic successor of $A$, so
$t(X)=\eta_+$ and \eqref{aspr:last-root} gives $k$. These are all
positions. Finally \eqref{aspr:near-free}, applied separately to the
parent and core, restores their geometric near arrays without changing
their complete root outputs. This proves the formal duplication identity.

\emph{Specialization to the soft polygon.}
Write $A=\mu_j$ and $B=A+\varepsilon q$, and let the parent have
$n+1$ vertices, so the fixed core $P$ has $n\geq3$ vertices. Every
parent determinant not containing both $A,B$ either is a core
determinant or replaces $A$ by $A+\varepsilon q$ in one. For example,
bilinearity gives
\begin{equation}\label{aspr:det-clearance}
\det(X-A-\varepsilon q,Y-A-\varepsilon q)
=\det(X-A,Y-A)-\varepsilon\det(q,Y-X).
\end{equation}
The potential quadratic term vanishes because $\det(q,q)=0$;
alternation gives the same bound when $A$ is in another position.
The core has finitely many nonzero distinct-triple determinants, so
their absolute values have a positive minimum $\delta$. The finitely
many coefficients $\det(q,Y-X)$ have a finite maximum absolute value
$K$. Choose a common positive tail with $\varepsilon K<\delta$ when
$K>0$; if $K=0$, no restriction is needed. On this tail every such
parent sign equals its pulled-back core sign. All core edges have a
positive minimum length; shrinking further if necessary also keeps
the return edge $\widetilde\lambda_j-\varepsilon q$ nonzero. Take this
tail inside $(0,\varepsilon_0)$.

For every other core vertex $X$, set
\begin{equation}\label{aspr:geometric-t}
t(X)=\chi(B,A,X)=-\operatorname{sgn}\det(q,X-A).
\end{equation}
Indeed $A-B=-\varepsilon q$ and $X-B=X-A-\varepsilon q$, so their
determinant is $-\varepsilon\det(q,X-A)$. Admissibility makes it
nonzero. For a boundary triple $(X,A,B)$ the far entry is $t(X)$;
for $(A,B,Y)$ it is $\chi(Y,B,A)=\chi(B,A,Y)=t(Y)$ by a cyclic
permutation. Together with the common-tail calculation, these are
exactly \eqref{aspr:far-data}.

The two cyclic neighboring values are the stated attachments. Writing
$u=\mu_j-\mu_{j-1}$ and $v=\mu_{j+1}-\mu_j$, the predecessor gives
\begin{equation}\label{aspr:attachment-minus}
t(\mu_{j-1})=-\operatorname{sgn}\det(q,-u)
             =-\operatorname{sgn}\det(u,q)=\chi_-;
\end{equation}
the equality follows from bilinearity and antisymmetry. The successor
gives directly
\begin{equation}\label{aspr:attachment-plus}
t(\mu_{j+1})=-\operatorname{sgn}\det(q,v)=\chi_+.
\end{equation}

Cut the boundary word immediately after the specified nonsoft root.
The occurrences $A,B$ are consecutive in this linear word. If the root
is the incoming edge ending at $A$, the word is soft-first ($s=0$).
If it is the return edge starting at $B$, the word is soft-last
($s=M$). Every other nonsoft root gives the strictly spanning case.
Here the core word has $M=n-1\geq2$, so the formal result applies even
to a triangle core; no two-gon amplitude is invoked. Collapse of $B$
to $A$ retains each other root's physical endpoints, except that the
return edge becomes the core edge $E_j=[\mu_j,\mu_{j+1}]$. Thus the
collapsed word is exactly the boundary word of $(P,g')$, including
cyclic wraparound. The soft root itself would put $A,B$ at opposite
ends of the cut word and has been excluded throughout.

The duplication identity now gives \eqref{aspr:soft-law} with
$k=(\chi_-+\chi_+)/2$. Evaluating this expression on the three
attachment-sign patterns gives the sector multipliers in the
statement. Neither the crossing classification of the soft family
nor a rotation formula for it was needed in this calculation.
\end{proof}

\begin{lemma}[the local vertex identity]\label{lem:softvertex}
Let $P$ be generic, $q$ admissible and $0<\varepsilon<\varepsilon_0$ as in
Lemma~\ref{lem:soft-generic}; in the soft insertion, let $\pi$ be a
composition of an interval of the boundary word of $(P_\varepsilon,g)$ whose
first part is the single soft leaf, $\pi=(\mu_j,\mu_*,a_2,\ldots,a_m)$ in the
main text's notation $(1,2,a_2,\ldots,n)$. Then, for every such $\varepsilon$,
\[
V^\pm(\mu_j,\mu_*,a_2,\ldots,a_m)=
\frac{\chi(\mu_*,\mu_j,a_2)\mp\chi(\mu_*,\mu_j,a_m)}{2}\,
V^\pm(\mu_j,a_2,\ldots,a_m),
\]
the unlabelled local vertex identity following the main text's
\mt{eq:softinsertionlaw}. \status{proved (refereed: bench A 2026-09-12T07:48:52Z, bench B 2026-09-12T07:39:15Z; SM1-carried, read on SM12 by both benches, repaired in round 12 (F-25-192) and round 13 (F-25-198), re-read on SM13 and SM14 by both benches)}
\end{lemma}
\begin{proof}
By Lemma~\ref{lem:soft-generic}(i) $P_\varepsilon$ satisfies (G1), so
Lemma~\ref{lem:gates-nonzero} applies to it: the gate at the cut $\mu_*$ has
near sign $d=\chi(a_2,\mu_*,\mu_j)=\chi(\mu_*,\mu_j,a_2)$ and far sign
$h=\chi(a_m,\mu_*,\mu_j)=\chi(\mu_*,\mu_j,a_m)$ (cyclic permutations), and its
ordinary and root values are $(d\mp h)/2$. The right side is read in the
boundary word of $(P,g)$, identified with that of $(P_\varepsilon,g)$ with
$\mu_*$ deleted through the identification of the edges of $P_\varepsilon$
with those of $P$ in Lemma~\ref{lem:soft-generic}(iii). Every other gate of
$\pi$ reads a triple in which $\mu_*$ occurs only as the neighbour $a_1$ of
the cut $a_2$, and for $0<\varepsilon<\varepsilon_0$ the determinant of the
triple $(a_3,a_2,\mu_*)$ is continuous in $\varepsilon$ and nonzero by (G1)
of $P_\varepsilon$ (Lemma~\ref{lem:soft-generic}(i)), hence of constant sign
on that connected interval, equal to its nonzero limit $\chi(a_3,a_2,\mu_j)$
as $\varepsilon\to0$; the far signs of $\pi$ do not involve $\mu_*$.
\end{proof}

\subsection{Reversal and shift}

\begin{proposition}[reversal and cyclic shift]\label{prop:A-reversal}
For every polygon satisfying (G1) and every root $g$:
\begin{enumerate}
\item[(i)] $A_g(\sigma P)=A_{g+1}(P)$.
\item[(ii)] With $\overline P_i=\mu_{2-i}$ and $g^\vee=1-g$,
\begin{equation}\label{arpr:reversal}
A_{g^\vee}(\overline P)=(-1)^nA_g(P).
\end{equation}
\end{enumerate}
All vertex and edge indices in this statement are read modulo $n$.
\status{proved (refereed: bench A, 2026-09-05T15:51:04Z; transcribed from RC mt:covariance; appendix\_v7 eq:treecuts (C017; F-25-20))}
\end{proposition}
\begin{proof}
The relabellings preserve (G1), because they permute distinct vertex
triples. Put $N=n-1$. For the shift, the boundary vertex at position
$k$ of $(\sigma P,g)$ is $\mu_{g+2+k}$, which is precisely position
$k$ of $(P,g+1)$. Every interval composition, near sign and far sign
is therefore identical. Induction on interval length in the open-sum
recursion gives equality of all open sums, and the complete root sums
agree, proving (i).

For reversal, edge $g^\vee$ runs from
$\mu_{2-g^\vee}=\mu_{g+1}$ to
$\mu_{1-g^\vee}=\mu_g$, the reverse of the original physical root.
If $a_k=\mu_{g+1+k}$ is the original boundary word and $a'_k$ the
word at this reversed root, then
\begin{equation}\label{arpr:boundary}
a'_k=\overline\mu_{g^\vee+1+k}
     =\mu_{1-g^\vee-k}=\mu_{g-k}=a_{N-k}
\qquad(0\leq k\leq N).
\end{equation}
In the last equality the indices differ by $n$.

Use the tree expression of Lemma~\ref{lem:treesum-trees}.
Reflect the order of the children at every internal vertex of a
rooted plane tree, and relabel boundary positions by $k\mapsto N-k$.
This is an involution on the tree set: the root is still allowed one
or more children, ordinary vertices still have at least two, and leaf
order is exactly the reversed boundary order. A composition
$(r_0<\cdots<r_s)$ becomes
$(N-r_s<\cdots<N-r_0)$ on the reflected interval. At corresponding
cuts its ordered near triple reverses its first and last entries;
the same is true of its ordered far triple. Thus $d'=-d$ and $h'=-h$.
The product $h'd'=hd$, so either gate retains its step-function
argument and acquires precisely one minus sign:
\begin{equation}\label{arpr:gate}
d'\Theta(\mp h'd')=-d\Theta(\mp hd).
\end{equation}
This holds for both ordinary and root gates, including zero factors.
The separate factor $(-1)$ for each ordinary internal vertex is
unchanged by the involution.

If the tree has $I$ internal vertices and $N$ leaves, its number of
child edges is $I+N-1$. The sum of the arities over internal vertices
counts each child edge once. A vertex of arity $s$ contributes $s-1$
gates, so the total number of gates is
\begin{equation}\label{arpr:gate-count}
\sum_{v\text{ internal}}(s_v-1)
=(I+N-1)-I=N-1=n-2.
\end{equation}
A unary root contributes zero to this count, so its allowed presence
creates no exception. Every corresponding tree term therefore differs
by $(-1)^{n-2}=(-1)^n$. Summing over the involution gives
\eqref{arpr:reversal}. For $n=3$ the same count is one; no exceptional
base convention is required. Neither part proves or assumes root
independence on a fixed polygon.
\end{proof}

\begin{remark}[rooted scope]\label{thm:root-independence}
The construction and the local laws in this section retain the specified
root. None of their proofs assumes root independence on a fixed polygon.
\end{remark}
